% Draft section authority: Paper/afmr_story_spec.md, src/eviseq_new,
% and Technical_Report/AFMR_IMPLEMENTATION_PLAN.md (Sections 1--10, 27--28).
% The surrounding manuscript should define the model, dataset, and metric
% commands used below in its preamble.

\section{Introduction}
\label{sec:introduction}

% Role: establish the task and the factuality--overlap distinction.
Long-document abstractive summarization turns a lengthy source into a shorter text. For scientific, biomedical, book, and government documents, a useful summary must preserve entities, numbers, and relations that the source supports. Factual support is therefore different from lexical overlap: a fluent summary can resemble its reference while changing a quantity or relation. Prior work has documented this gap in abstractive summarization \citep{maynez2020faithfulness}.

% Role: state the interface gap created by independently pretrained components.
Most sequence-to-sequence models learn source encoding and target generation through one interface, with the decoder attending to encoder states \citep{sutskever2014sequence,bahdanau2015neural,vaswani2017attention}. T5, BART, and mT5 pretrain this interface jointly \citep{raffel2020t5,lewis2020bart,xue2021mt5}. A different problem appears when an embedding-oriented encoder and a causal decoder are reused from separate checkpoints. Their hidden widths, representation coordinates, tokenizers, and pretraining objectives may all differ. A projection can make the widths compatible, but it does not determine which source representations, positions, or lexical forms should guide each generated token.

% Role: sharpen the technical challenge and bridge to the proposed interface.
This composition creates two related problems. The interface must keep every visible source position available while assigning useful priority across encoder depth, feature channels, and source spans. It must also expose source words and other lexical forms without allowing a separate retrieval objective to dictate generation. Long-document attention and copying address parts of these problems, but they do not specify how to allocate information between independently pretrained representation spaces \citep{cohan2018discourse,see2017get}. We frame that allocation as an interface problem and test it with the same generation graph and matched comparisons.

% Role: give the method idea at the level needed to understand the questions.
We study this problem with \EviSeq and its proposed \AFMRFull{} (\AFMR) interface. \AFMR keeps a full retrieval tensor for every visible source position alongside an aligned final-state value tensor, then learns small adjustments conditioned on the source and decoder budgets. These adjustments act along three axes. A depth route combines the final encoder state with the last four hidden-state taps around a final-state anchor. A cross-space feature route adapts the encoder representation to the decoder width while leaving the main projection path available. A multi-scale route builds a shared additive prior from overlapping source-token windows of lengths 32, 128, and 512 in the base profile. Every decoder layer reads these tensors through a copied cross-attention block: refined representations form the keys, and final-state representations supply the values. A grounded-copy route aligns source and decoder-token offsets, combines lexical and contextual keys, and marginalizes repeated source token IDs. The model keeps all visible source positions and needs no evidence labels at inference.

% Role: define the evidence boundary and the locked comparison protocol.
The graph uses the gold-token mixture cross-entropy that defines generation. Training begins with one interface warm-up epoch and then moves to full fine-tuning; official generation uses one beam with sampling disabled. We apply the same preprocessing, splits, source and target budgets, and training budgets to \PubMed, \ArXiv, \BookSum, and \GovReport. The comparisons include fine-tuned decoder-only controls and fine-tuned \TfiveGemma. Implementation and unit tests can verify tensor, masking, gradient, checkpoint, and inference invariants, but the four-dataset predictions and quality measurements are not available at this stage. \AFMR is consequently a testable hypothesis, and we make no claim of improved \Rouge or factuality before the registered runs.

% Role: pose the locked four-question chain.
Across \PubMed, \ArXiv, \BookSum, and \GovReport, we ask four linked questions:

\begin{itemize}
  \item \textbf{RQ1 (factual-support signal):} Does \EviSeq reduce a source-unsupported-content signal relative to fine-tuned decoder-only LLMs on all four datasets? We use native \AlignScore consistency and its direction-transformed $H=1-C$ value for the automatic comparison; a claim about typed entity, number, or relation errors requires a separate adjudicated audit.
  \item \textbf{RQ2 (quality):} How does \EviSeq change \Rouge-1/2/L and \BERTScore relative to the same decoder-only controls and fine-tuned \TfiveGemma under matched records, budgets, preprocessing, and decoding? Each model's native input serialization is documented for all four datasets.
  \item \textbf{RQ3 (mechanism):} Which of the encoder, \AFMR bridge, and grounded-copy route contributes to the observed behavior? We run no-encoder, no-bridge, and no-grounded-copy controls on all four datasets with a fixed decoder and matched training budget.
  \item \textbf{RQ4 (portability):} Does the same interface transfer across a bidirectional embedding encoder, a causal encoder, and a larger embedding encoder, and how do directionality and capacity affect quality and factuality? We use \PPLXEmbed as the anchor and test \QwenCausal, \QwenEmbed, and \NemotronEmbed variants on all four datasets.
\end{itemize}

% Role: make the causal bridges between the questions explicit.
Together, the four questions separate source support, summary quality, component contribution, and encoder portability. RQ1 tests whether the interface changes an automatic source-support signal. RQ2 checks whether that change comes with different reference overlap or semantic similarity. RQ3 identifies the routes that contribute to the observed benchmark behavior, while RQ4 asks whether the behavior persists across encoders with different directionality and capacity. For RQ1, higher native \AlignScore consistency is better. We report $H=1-C$, where $C=\AlignScore_{\texttt{nli\_sp}}$, only as a direction-transformed hallucination-oriented score; lower is better, and $H$ is not a calibrated probability. RQ2 uses \Rouge-1/2/L F1 and \BERTScore-F1 as complementary measures of overlap and semantic similarity. All interpretations are conditional on the held-out predictions and paired uncertainty intervals.

% Role: state contributions without converting planned evidence into results.
This paper contributes three things. First, it formulates \AFMR as a constrained full-memory interface for composing a pretrained encoder with a causal decoder, with key-side adaptation separated from final-state values. Second, it combines controller-conditioned depth, cross-space feature, and multi-scale span adjustments with a grounded-copy route while keeping one gold-token generation objective. Third, it defines a reproducible four-dataset evaluation contract that separates source-support signal, overlap quality, component contribution, and encoder portability, and records the diagnostics and uncertainty estimates needed to interpret them. The interface is experimentally falsifiable; these contributions do not imply that any component is necessary or that the complete system is superior before the planned comparisons.

\section{Related Work}
\label{sec:related-work}

% Role: position AFMR against jointly trained and composed encoder--decoder systems.
\paragraph{Encoder--decoder interfaces and checkpoint composition.}
Sequence-to-sequence learning and attention established the division between source encoding and autoregressive generation \citep{sutskever2014sequence,bahdanau2015neural,vaswani2017attention}. T5, BART, and mT5 learn this interface during joint pretraining, so their source and target representations are optimized together \citep{raffel2020t5,lewis2020bart,xue2021mt5}. Other work reuses decoder-only checkpoints for representation learning, including LLM2Vec and Generative Representational Instruction Tuning \citep{behnamghader2024llm2vec,muennighoff2024grit}. LaMaTE uses a language-model encoder in a translation system, while Qwen3 Embedding provides an embedding-oriented model family \citep{luo2025lamate,zhang2025qwen3embedding}. \TfiveGemma is a close reference for arranging modern language-model components as an encoder--decoder system \citep{zhang2025t5gemma2}. \AFMR shares the goal of reusing pretrained components and the decoder's cross-attention interface, but addresses the allocation problem that remains after composing checkpoints that were not released as a matched encoder--decoder model. We claim no novelty for generic checkpoint reuse or cross-attention itself.

% Role: connect long-document methods to the full-memory design distinction.
\paragraph{Long-document source access.}
Long-document summarization has used hierarchical and discourse-aware attention to represent documents at several granularities \citep{cohan2018discourse}. These methods motivate looking at both local and broader source regions, but they do not decide how a composed model should turn regional preference into decoder retrieval. \AFMR uses overlapping token windows at several scales to build a soft source prior. The decoder still receives retrieval and aligned value tensors for every visible source token. This leaves hard selection, sentence slots, and source compression outside the proposed interface and allows direct comparison with no-bridge and matched-encoder controls.

% Role: distinguish grounded copying from pointer and coverage mechanisms.
\paragraph{Copying and source coverage.}
Pointer-generator networks mix vocabulary generation with an attention-based copy distribution, while coverage mechanisms can discourage repeated attention to the same source content \citep{see2017get}. \AFMR uses this generation--copy decomposition for independently tokenized components. Character-offset overlap supplies a sparse source-to-decoder alignment; contextual and lexical features form the copy keys; and repeated source occurrences are summed by decoder token ID. The copy route is trained with the same token-level mixture likelihood as the language-model path. \AFMR has no separate coverage loss and does not treat copying as factual verification: a copied source token can still appear in the wrong relation or role.

% Role: separate factuality evidence from lexical and semantic similarity metrics.
\paragraph{Factuality and summary evaluation.}
\Rouge measures overlap with a reference \citep{lin2004rouge}, whereas \BERTScore measures contextual semantic similarity between generated and reference texts \citep{zhang2020bertscore}. Neither metric directly guarantees factuality, as prior work on faithfulness and evaluation has shown \citep{maynez2020faithfulness,fabbri2021summeval}. We therefore use \AlignScore as the primary automatic source-consistency signal, report its native direction, and show $H=1-C$ for $C=\AlignScore_{\texttt{nli\_sp}}$ only as a lower-is-better diagnostic \citep{zha-etal-2023-alignscore}. \AlignScore is still a learned proxy and may miss cross-sentence or domain-specific errors; \BERTScore is not a hallucination detector. Our evaluation pairs these metrics with prediction diagnostics and, where available, a manual factual-error audit.

% Role: close Related Work with the precise distinction and evidence boundary.
These lines of work provide the pieces used by \AFMR: source--target attention, multi-scale source reasoning, and vocabulary--copy mixtures. \AFMR brings them together at a full-memory boundary, adds bounded controller-conditioned residuals for encoder depth, cross-space features, and source-span priority, and trains the graph with gold-token cross-entropy. RQ3 tests whether these additions matter with a fixed decoder and matched budget; RQ4 tests transfer across the specified encoder families. The literature motivates the design, but only the locked four-dataset experiments can support claims about quality, factual support, component contribution, or portability.
